We consider Federated Learning constructed over $n = 1, \ldots, N$ communication rounds, each consisting of $k = 1, \ldots, K$ local update steps.
The set $C = \left[1, \ldots, M \right]$ represents the users participating in federated learning, each with local objective function $F_i$. Federated learning aims to find the model that minimizes the average user objective.
Note we distinguish $F_i$, which is evaluated on the entire local data set, from $F_i^{(n, k)}$, which is evaluated on data from the $k$-th local update step in round $n$.
In each round, a subset of users, $S \subseteq C$, participates in the update of the global model, and the case $S = C$ is full user participation.
We denote the global model at communication round $n$ as $\theta_g^{(n)}$, and $\theta_i^{(n, k)}$ represents the $i$-th user's copy of the shared model after $k$ of $K$ planned local update steps with learning rate $\eta$.
This section begins with introducing our notation for the federated learning process.

Before introducing the proposed method, we provide background information on the similarities and differences between the popular existing methods: FedSGD, FedAvg \citep{fedavg}, and FOMAML \citep{reptile}. 
We believe this viewpoint of existing algorithms is essential for understanding the derivation and benefits of our proposal.
Importantly, we justify that our proposed method can flexibly balance the goals of \textbf{initial model success}, shared model performance on the local datasets, and \textbf{rapid personalization}, personalized model performance on local data after fine-tuning.
Mathematically, we define initial model success in Equation~\ref{eq: ims} as the average user objective with the shared model at the end of federated training. 
Similarly, rapid personalization in Equation~\ref{eq: rp} is the average user objective with the personalized model, fine-tuned with the local data, after federated training is complete as indicated by the superscript $(N + 1)$. 
\begin{align}
        F_{IMS} &= \frac{1}{M} \sum_{i = 1}^M F_i (\theta_g^{(N)}) \label{eq: ims} \\
    F_{RP} &= \frac{1}{M} \sum_{i = 1}^M F_i(\theta_i^{(N + 1)}) \label{eq: rp}
\end{align}

We remark that either of these objectives could be weighted.
However, in this work, we focus on an algorithm that performs well for all users and treats all users equally. 
Note that $F_{IMS}$ is the same objective as commonly minimized in traditional Federated Learning and $F_{RP}$ for Personalized Federated Learning (and meta-learning). 
We will demonstrate that controlling the application-specific focus between initial model success and rapid personalization improves personalization and can extend to new users without sufficient data for fine-tuning.
Contrary to intuition, we find that relative to previous applications of meta-learning, placing additional emphasis on initial model success can improve performance after rapid personalization. 
We believe this applies to most federated learning scenarios where some similarities exist between users. 
Furthermore, we demonstrate how to maintain the same rate of convergence as FedAvg under certain assumptions while preserving the interpretability of the model performance trade-off mentioned above. 
